\section{模型假设}

为统一小时级调度与分钟级资源计量，假设任务开工时刻按照整数小时选取，而任务在各小时内的 GPU-hour 和能耗根据实际起止时刻与该小时的重叠时长计算。附件中的 GPU 需求、预计持续时间、到达时刻、来源区域、最大允许时延和最晚完成时刻均被视为准确且确定的输入参数，不另行建立执行时间和资源需求的随机模型。所有任务仅能选择一个执行区域和一个开工时刻，自开工至完成保持连续运行，不允许拆分、抢占、暂停或中途迁移。

假设实时推理任务必须在到达时刻立即开工，其执行区域只能从满足单向网络时延和资源容量的区域中选择。批量推理与 AI 训练任务可以延迟开工，但不得早于到达时刻，且必须在任务有效截止时点与系统时点 $2406$ 中较早者之前完成。第 $2400$ 至第 $2405$ 小时只用于此前到达弹性任务的结清，第 $2406$ 小时不产生任务资源占用，只承担电力与储能终端结算。

假设同一任务类型的单位等效 GPU 平均 IT 功率在研究期内保持为功率映射附件给定值，不随区域和时间改变。非 AI IT 负荷在各区域和各时段固定且不可迁移，只有 AI 任务负荷能够随执行区域及开工时刻改变。区域设施功率由 AI 与非 AI 合计 IT 功率乘以对应 PUE 得到，PUE 在研究期内按附件取值，不建立其随环境或负荷率变化的动态模型。

假设区域间通信只体现为附件给出的来源区域至执行区域单向网络时延，时延在研究期内确定且不受网络拥塞影响。来源区域与执行区域相同时仍采用时延矩阵中的对应记录，不自行设定为零。由于题面没有给出网络带宽、迁移数据量、传输能耗和传输费用等信息，上述因素不被纳入模型。

假设购电价格、售电价格、电网碳强度、可用新能源和固定负荷在单个小时内保持恒定。各区域逐小时独立满足能源平衡，不建立区域间线路潮流模型，区域售电被视为向系统边界外送。可用新能源可以被分配至直接供给负荷、储能充电、外送或弃电，其优先关系由实际策略与能源平衡共同确定，而不预设全部新能源均被本地直接消纳。运行碳排放仅按电网购电量与相应碳强度累计，不计新能源、储能设备和通信设施的生命周期排放。

假设储能在第 $0$ 小时运行前的状态唯一取自附件中的初始 SOC，后续状态均在该绝对值基础上递推，不将初始状态归零或作为待估参数。区域逐时基准 SOC 只用于方案核验与比较，不作为优化轨迹的固定约束。充电效率作用于进入储能的能量，放电效率作用于储能侧能量消耗，未给出的自放电与储能退化成本不予考虑。每个区域在第 $2406$ 小时运行后的 SOC 不低于其初始 SOC，从而避免通过透支终端能量获得虚假的当期收益。

假设同一储能在同一小时只能处于充电、放电或空闲状态，不能同时充放电；区域购电和售电也不能在同一小时同时发生。两类行为均受附件所给功率上限约束，若价格和效率条件不能自然排除无效循环，则通过互斥状态约束加以限制。问题三中的 IT 负荷固定为基准 AI IT 负荷与非 AI IT 负荷之和，不重新调整任务的执行区域和开工时刻。

假设预测模型与正式调度严格分离，预测区间的估计值只用于测试精度评价，最后 $24$ 小时基础调度以及问题二至问题四的正式计算均使用实际任务和实际逐时参数。服务质量仅由题内可观测的即时开工、单向时延 SLA、按期完成情况及弹性任务等待时间表示，不引入用户满意度或可靠性等缺乏数据支持的变量。进行碳约束、电价机制与新能源波动情景比较时，除指定变化因素外，其余任务集、设备容量、网络时延、功率映射及评价公式均保持一致。

假设普通空白值在未经核验前均属于不可裁决缺失，不自动替换为零。只有在互斥进流与出流字段中，空白所在行为孤立事件、对应另一列明确非零，且备注、相邻记录和全表口径共同支持互斥解释时，才将其认定为结构零。经过处理的数据还应接受任务唯一执行、实时即时开工、时延与截止、GPU 和功率容量、能源守恒、SOC 递推、终端状态及购售电边界的统一审计。
